\chapter{Sprint 43 --- Công nợ \& thanh toán portal mobile, 7 màn Figma (2026-07-31)}

\section{Bối cảnh}

Cụm Figma \texttt{WJ\_Debt\_*\_MVP\_v31} (page \emph{Dashboard}, node \texttt{5013:*}) gồm 7 frame
391$\times$844 là cụm thiết kế \textbf{tồn cuối cùng chưa có code}. Đồng thời portal đang mang hai
\textbf{điểm vào chết}: tile KPI \emph{Công nợ} ở Home mobile là ô \texttt{0 đ} trơ với comment
\emph{``module /portal/debt đã xoá''}, và mục \emph{Công nợ} ở sidenav PC trỏ ngược về
\texttt{/portal}. Sprint này dựng lại trang để hai điểm vào đó có đích thật.

\section{Vì sao UI-only}

Quyết định này không phải để làm nhanh; bốn dữ kiện độc lập đều chỉ về một hướng:

\begin{itemize}
  \item Tab \texttt{1. Model/ Field} mục \textbf{``D. Quản lý công nợ nhượng quyền'' (dòng 529) là
    tiêu đề rỗng} --- BA chưa đặc tả \emph{một} model/field nào. Tab \texttt{3. Controller} có
    CT-050\ldots CT-055 nhưng cột ghi chú còn để ngỏ (\emph{``cần chốt QR tĩnh hay QR động''}), và
    tab \texttt{Tasks} không có dòng nào cho công nợ.
  \item Database: \texttt{account} / \texttt{account\_payment} đã cài nhưng có đúng \textbf{1
    \texttt{account.move}} (bút toán mở sổ), \textbf{0 \texttt{account.payment}}, và \textbf{không
    field nào nối \texttt{account.move} với cửa hàng nhượng quyền}. Không có gì để đọc.
  \item Chính Figma tự ghi \emph{``QR minh họa $\bullet$ không dùng để thanh toán''} và
    \emph{``Vietcombank (minh họa)''}.
  \item Tiền lệ: module \texttt{wujia\_portal\_debt} bản cũ (dựng Sprint 16, xoá Sprint 17) có
    docstring ghi rõ \emph{``card Công nợ = UI-ONLY, backend account.move thuộc Phase 2''}.
\end{itemize}

Nói cách khác, wire backend bây giờ là \textbf{tự bịa schema tài chính} rồi bắt BA sống với nó.

\section{Kiến trúc --- một seam dữ liệu, không sửa file module cũ}

Module mới \texttt{wujia\_portal\_debt} (tên tái dùng được: bản ghi
\texttt{ir\_module\_module} đang \texttt{uninstalled}), 13 file / 1529 dòng, \texttt{depends:
['wujia\_portal\_base']}.

\subsection{Nguồn dữ liệu duy nhất}

\texttt{models/wujia\_portal\_debt.py} khai báo \texttt{models.AbstractModel} tên
\texttt{wujia.portal.debt} --- \textbf{không bảng, không migration, không có gì để gỡ}. Bốn method
\texttt{get\_summary} / \texttt{get\_payments} / \texttt{get\_bank\_info} / \texttt{get\_shell\_badge}
trả dict thuần Python (0 query). Controller, template và cả hai view inherit đều \emph{chỉ} gọi qua
đây. Khi BA chốt spec backend, việc phải làm là đổi nguồn bên trong bốn hàm đó; template không đụng
một dòng.

Test \texttt{test\_shell\_badge\_shape} chốt cứng \textbf{shape} của dict badge --- đổi shape là vỡ
Home và sheet \emph{Thêm}, và test nói ra điều đó ngay thay vì để lộ trên UAT.

\subsection{Sửa shell bằng inheritance, không sửa module shell}

Hai thay đổi lên vỏ portal đều làm \textbf{từ module mới}:

\begin{center}
\begin{tabular}{p{4.9cm}p{8.9cm}}
\hline
\textbf{View} & \textbf{Việc} \\
\hline
\texttt{bottomnav\_inherit.xml} & \texttt{inherit\_id="wujia\_portal\_layout.mobile\_bottomnav"},
xpath chèn dòng \emph{Công nợ \& thanh toán} + badge lên đầu \texttt{.wujia-msheet-list}. \\
\texttt{home\_kpi\_inherit.xml} & \texttt{inherit\_id="wujia\_portal\_base.portal\_home\_page"},
thay tile KPI trơ bằng \texttt{<a href="/portal/debt">} + số còn phải trả thật. \\
\hline
\end{tabular}
\end{center}

Hệ quả: \texttt{wujia\_portal\_layout} và \texttt{wujia\_portal\_base} \textbf{không sửa một dòng
nào}, nên blast radius đúng bằng việc cài/gỡ \texttt{wujia\_portal\_debt}.

\begin{quote}
\textbf{Ràng buộc ghi lại cho lúc wire thật.} Sheet \emph{Thêm} nằm trong shell $\Rightarrow$
\textbf{mọi trang mobile} đều gọi \texttt{get\_shell\_badge()}. Hiện là dict thuần nên 0 query;
lúc nối dữ liệu thật, số ``quá hạn'' \textbf{bắt buộc} là field store+index hoặc
ormcache/cron daily --- \emph{không} \texttt{search\_count} mỗi request (1500 portal user, §7).
\end{quote}

\subsection{Route}

\texttt{GET /portal/debt} (4 biến thể theo \texttt{state}, query \texttt{week} + \texttt{all=1}) $\bullet$
\texttt{GET /portal/debt/payment-history} $\bullet$ \texttt{GET /portal/debt/pay}. Cả ba dùng
\texttt{get\_active\_franchise\_id()}; chưa chọn cửa hàng thì \textbf{render empty state, không 500}.

\section{Dựng 7 màn --- và giữ nguyên chỗ Figma tự mâu thuẫn}

Khung chung: content rộng \texttt{359}, padding ngang \texttt{16}, card trắng radius 16, nền
\texttt{--wujia-bg-page}.

Điểm đáng ghi: card tổng ở frame 02 xếp stats thành \textbf{2 cột}, còn frame 03/04 xếp \textbf{2
dòng} --- vì frame 02 có thêm dòng ``hạn gần nhất'' mà cả ba card đều cao đúng 142. Đây là chỗ dễ
``tự chuẩn hoá cho đẹp''. \textbf{Không chuẩn hoá}: dựng đúng như Figma bằng modifier class theo
state, và ghi lại lý do ở đây để lần sau không ai tưởng là bug.

Frame 05 (empty) cũng vậy: card rỗng \emph{vẫn} giữ row \emph{Lịch sử thanh toán} bên dưới --- không
có công nợ tuần này không có nghĩa là không có lịch sử để xem.

\subsection{Màu --- một sai lệch có chủ đích}

Mọi hex trong Figma đều đã có token trong \texttt{\_variables.css}, nên CSS module chỉ khai báo alias
cục bộ \texttt{--wj-debt-*: var(--wujia-*)}. Còn \textbf{đúng một} hex cứng:
\texttt{--wj-debt-warn-strong: \#B45309} cho chữ tiêu đề trên nền cảnh báo \texttt{\#FFF7E6}
(token \texttt{--wujia-mres-warn-text} \texttt{\#D97706} không đạt AA trên nền đó) --- đề xuất BA
nâng lên token global.

Sai lệch cần BA xác nhận: Figma vẽ CTA \emph{Thanh toán} nền \texttt{\#28A9DF} chữ trắng. Trắng trên
\texttt{\#28A9DF} \textbf{trượt WCAG AA}, và Sprint 38 đã đổi CTA portal sang \texttt{--wujia-cta}
\texttt{\#0F7CA8} đúng theo issue a11y do chính BA mở. Dùng \texttt{--wujia-cta}.

\section{Đối chiếu Figma bằng số đo, không bằng mắt}

Theo lesson L4, việc đối chiếu được viết thành \textbf{harness đo được} thay vì nhìn ảnh: script
Playwright duyệt 7 màn ở 391$\times$844, so \texttt{bounding-box} / \texttt{computed-style} với số
lấy thẳng từ node JSON Figma, cộng hai kiểm tra tự động cho mỗi trang --- \emph{nội dung có tràn card
height cứng không} và \emph{khoảng cách giữa các khối có đúng không}.

Harness bắt được năm lỗi mà mắt thường bỏ qua, và cả năm đều là \textbf{shell Vuexy đè component}:

\begin{center}
\begin{tabular}{p{4.4cm}p{9.4cm}}
\hline
\textbf{Triệu chứng đo được} & \textbf{Nguyên nhân thật} \\
\hline
Title \texttt{client=278 scroll=283} (ellipsis) & \texttt{body \{letter-spacing: 0.14px\}} toàn cục
$\Rightarrow$ tiêu đề 26px tràn. \\
Card tổng: nội dung 150 trong hộp 142, đè lên CTA & \texttt{line-height: 1.8} toàn cục $\Rightarrow$
phải khai báo lại leading Figma 15/15/18. \\
Gap filter\,--\,card = 27 thay vì 12 & Filter là thẻ \texttt{<form>}, Vuexy có
\texttt{form \{margin-bottom: 15px\}}. \\
Dòng sheet \emph{Thêm} cao 83 giữa 6 dòng 65 & Shell chạy type-scale 15/12 còn Figma vẽ 14/11 nên
badge ở cấp hàng đẩy phụ đề xuống 2 dòng. Chuyển badge vào \emph{trong} dòng title $\to$ 68px, phụ đề
1 dòng, chữ Figma giữ nguyên. \\
Sửa CSS xong không thấy đổi gì & Asset trong \texttt{web.assets\_frontend} được cache ở
\texttt{ir.attachment} $\Rightarrow$ phải \texttt{-u} hoặc chạy \texttt{-{}-dev=xml,assets}. \\
\hline
\end{tabular}
\end{center}

Harness cũng \textbf{tự tố cáo hai báo động giả} của chính nó, và đó là phần có giá trị nhất: nút
back đo 40 trong khi Figma ghi 42 (chênh do \texttt{border-box}), và gap đầu tiên ở hai trang có nút
back là 8 chứ không phải 12. Mở lại node JSON thì header kết thúc ở \texttt{y=202} còn khối kế bắt
đầu ở \texttt{210} --- \textbf{code đã đúng, kỳ vọng trong script mới sai}. Sửa kỳ vọng, không sửa
code. Badge tràn khỏi vùng \texttt{\_\_head} cao 15 cũng là \emph{cố ý} theo Figma nên được khai báo
miễn trừ tường minh trong harness thay vì nới lỏng ngưỡng cho cả trang.

\section{Kết quả kiểm thử}

\begin{itemize}
  \item \texttt{-u wujia\_portal\_debt -{}-test-enable -{}-test-tags wujia\_debt}: \textbf{RC=0,
    0 failed / 0 error trên 14 test}; \texttt{grep} ERROR/CRITICAL/Traceback trong log $=$ \textbf{0}.
  \item Harness đối chiếu Figma: \textbf{ALL OK} trên cả 7 màn --- content 359, card tổng 142,
    CTA 44, row lịch sử 54, card hoá đơn 62, card thanh toán 96, card QR 220 (khung 160 / mã 144),
    card bank 150, card số tiền 76, badge 26, title 26/22px, gap 12 (8 ở trang có nút back);
    không màn nào có nội dung tràn card.
  \item Smoke shell: \texttt{/portal}, \texttt{/portal/order}, \texttt{/portal/notification} ở cả
    391$\times$844 lẫn 1920$\times$1080 đều HTTP 200, \textbf{overflow ngang $=$ 0}; KPI Home nay là
    link thật đọc \emph{12,7tr} thay vì ô trơ \texttt{0 đ}. Console chỉ còn 404 sẵn có của Vuexy
    (\texttt{app-assets/data/locales/en.json}).
  \item Chụp 7 ảnh local đặt cạnh 7 PNG Figma: khớp; khác biệt còn lại thuần dữ liệu (tên/mã cửa
    hàng, ngày tuần demo, hoạ tiết QR placeholder).
\end{itemize}

\section{Chưa làm --- và vì sao}

Sidenav PC \textbf{không đụng}. Figma chưa có frame PC cho công nợ; trỏ mục PC sang
\texttt{/portal/debt} sẽ hiện layout mobile trong khung 1920 --- đúng loại lỗi
\texttt{UI-01}/\texttt{UI-PC-SHELL-001} mà chính BA đang bắt. Trang vẫn render an toàn ở desktop
(cột 391 căn giữa) để ai gõ thẳng URL không gặp trang vỡ, nhưng \textbf{chờ BA cấp Figma PC} rồi mới
mở điểm vào.

Sáu dòng cũ của sheet \emph{Thêm} giữ nguyên icon feather: Figma vẽ icon là chữ viết tắt
(\emph{CN/ĐT/ĐK}\ldots) nhưng đổi cả sáu dòng nằm ngoài phạm vi sprint.

Chưa deploy UAT --- module mới nên \textbf{bắt buộc \texttt{-i}} một lần, gộp vào đợt deploy tồn
S39/S40/S42.

\section{Bài học}

\begin{itemize}
  \item \textbf{Tiêu đề rỗng trong spec là một câu trả lời, không phải chỗ trống để điền.} Mục D có
    tiêu đề mà không có field nào bên dưới nghĩa là BA chưa quyết --- dựng UI trước, chừa seam, để BA
    quyết sau.
  \item \textbf{Harness sai thì sửa harness, đừng sửa code.} Hai trong bảy phát hiện là báo động giả;
    nếu tin script vô điều kiện thì đã sửa hỏng phần vốn đã khớp Figma. Nguồn chân lý là node JSON,
    không phải kỳ vọng do mình gõ ra.
  \item \textbf{Giữ nguyên chỗ Figma không nhất quán, và ghi lý do.} Stats 2 cột ở frame 02 / 2 dòng
    ở 03--04 trông như lỗi thiết kế nhưng là hệ quả của việc cả ba card đều cao 142.
  \item \textbf{Số đo lệch hầu như luôn là shell đè, không phải CSS của mình sai.} Cả năm lỗi thật
    lần này đều là rule tag-level toàn cục (\texttt{letter-spacing}, \texttt{line-height},
    \texttt{form margin}) --- trung hoà trong scope component, không sửa file shared.
  \item \textbf{Sửa CSS mà không thấy gì đổi thì nghi cache asset trước khi nghi selector.} Bundle
    \texttt{web.assets\_frontend} nằm trong \texttt{ir.attachment}; \texttt{-{}-dev=xml,assets} lúc
    lặp, \texttt{-u} sạch một lần lúc chốt.
\end{itemize}
